% ======================================================================
\section{Methodology}\label{sec:methodology}
% ======================================================================

This chapter describes the slicing algorithm as a sequence of
computational stages. Each stage is first motivated by the physical or
robotic requirement it serves, then formalized, and finally discussed
in terms of design decisions and edge cases. The Python realization of
each stage is described in Chapter~\ref{sec:implementation}.

Figure~\ref{fig:pipeline} gives the global picture. The pipeline has two
branches that share the same downstream machinery:

\begin{itemize}
  \item the \textbf{piece branch}: the structural part itself, sliced
        into non-planar layers with fiber-aligned infill;
  \item the \textbf{mold branch}: an automatically generated support
        structure printed before (underneath) the part, sliced into
        near-planar layers with a grid infill.
\end{itemize}

\begin{figure}[htbp]
\centering
\fbox{\parbox{0.92\textwidth}{\centering\vspace{2em}
\textbf{Input:} APDL model (solid mesh + mold surface) + fiber paths (CSV)\\[0.6em]
$\downarrow$\\
\textbf{Stage 1:} model orientation (mean-normal alignment with $+Z$) and lifting\\[0.4em]
$\downarrow$ \hspace{5em} $\downarrow$\\
\textbf{Stage 2:} piece slicing \hspace{2em} \textbf{Stage 2':} shadow volume $\to$ mold slicing\\
(isosurfaces of $q$) \hspace{2.6em} (cone / vertical projection, then $q$-slicing)\\[0.4em]
$\downarrow$ \hspace{5em} $\downarrow$\\
\textbf{Stage 3:} shells + infill per layer (dist-to-boundary field; fiber-aligned planes / fixed XY grid)\\[0.4em]
$\downarrow$\\
\textbf{Stage 4:} path stitching (greedy nearest-endpoint, bounded connections, blockers)\\[0.4em]
$\downarrow$\\
\textbf{Stage 5:} tool-pose computation (normal + tangent $\to$ orthonormal frame)\\[0.4em]
$\downarrow$\\
\textbf{Stage 6:} simplification (RDP / infill collapse) and multi-layer text export\\[0.6em]
\textbf{Output:} trajectory files for RoboDK
\vspace{2em}}}
\caption{Global structure of the slicing pipeline.}
\label{fig:pipeline}
\end{figure}

% ----------------------------------------------------------------------
\subsection{Model preparation: orientation and lifting}\label{sec:meth-orient}
% ----------------------------------------------------------------------

The printing process assumes a dominant build direction: the robot
approaches the part from "above" ($+Z$), and the mold is grown
"downward" from the part's underside onto the printing plane
$Z = z_{\mathrm{plane}}$. A structural model, however, arrives in an
arbitrary orientation in the global frame of the finite-element model.

\subsubsection{Mean-normal estimation and rotation}

The first step therefore estimates the model's natural build direction
and aligns it with $+Z$. The natural build direction is taken to be the
mean of the normals of all layer surfaces of the part. Writing
$\mathcal{S}_k$ for the $k$-th layer surface (their computation is
described in \S\ref{sec:meth-slicing}) and $\mathbf{n}_{k,j}$ for the
unit normal of cell $j$ of $\mathcal{S}_k$, the mean normal is the
surface-area-weighted-free average
\begin{equation}
  \bar{\mathbf{n}} \;=\; \frac{1}{N_{\mathrm{cells}}}
  \sum_{k} \sum_{j} \mathbf{n}_{k,j},
  \qquad N_{\mathrm{cells}} = \sum_k |\mathcal{S}_k|.
\end{equation}
The slicer then builds the rotation matrix $R \in SO(3)$ that maps
$-\bar{\mathbf{n}}$ onto $\mathbf{e}_z = (0,0,1)$:
\begin{equation}
  R = \exp\!\big(\theta\, [\mathbf{a}]_\times\big),
  \qquad
  \mathbf{a} = \frac{\bar{\mathbf{n}} \times \mathbf{e}_z}{\|\bar{\mathbf{n}} \times \mathbf{e}_z\|},
  \quad
  \theta = \arccos\!\big(\bar{\mathbf{n}} \cdot \mathbf{e}_z\big),
\end{equation}
with the degenerate case $\bar{\mathbf{n}} \parallel \mathbf{e}_z$
resolved by $R = I$. The rotation is applied \emph{in place} to all
mesh data of the model: the solid-mesh nodes, the mold mesh, its
Kd-tree, and its PolyData representation, so that every downstream
consumer sees a consistently rotated model.

\subsubsection{Lifting}

After rotation, the model is translated along $+Z$ by an offset so
that its lowest printed surface lies above the printing plane,
leaving room for the support structure. This is a pure translation of
the same data structures, again applied in place.

% ----------------------------------------------------------------------
\subsection{Piece slicing: layers as isosurfaces}\label{sec:meth-slicing}
% ----------------------------------------------------------------------

\subsubsection{Why isosurfaces}

A planar slicer intersects the solid with planes
$z = z_0 + k h$. The generalization adopted here replaces the planes
with the level sets of a scalar field $q$ defined over the volumetric
mesh. The field is chosen such that $q$ is (proportional to) the
signed distance along the local build direction, so that consecutive
integer values of $q$ are one layer height $h$ apart:
\begin{equation}
  \mathcal{S}_k \;=\; \{\,\mathbf{p} \in \mathcal{V} \;:\; q(\mathbf{p}) = q_k\,\},
  \qquad q_{k+1} - q_k = h \ \text{(in physical units)}.
\end{equation}
For a thin-walled part, the isosurfaces of $q$ follow the wall's
mid-surface: the layers are \emph{curved}, and the nozzle deposits
material conformally to the wall. The scalar field itself is produced
upstream by the \texttt{oamc} layer-initialization stage
(\texttt{model.init\_q()}); the slicer only consumes the field named
\texttt{"q"} stored on the volumetric grid, and extracts the level
surfaces with the VTK contour filter.

\subsubsection{Practical extraction}

Each layer surface is extracted with a single contour call on the
unstructured grid, at the level $q_k$. Two properties of the extracted
surface matter for everything downstream:

\begin{itemize}
  \item \textbf{Boundaries.} The isosurface of an open wall is an open
        surface: it has boundary edges (where the contour leaves the
        volume's boundary). These boundary curves are the reference for
        all shell offsets (\S\ref{sec:meth-shells}).
  \item \textbf{Duplicate points.} The contour filter can emit
        duplicate points along the surface boundary, which silently
        corrupts feature-edge detection. The slicer therefore merges
        coincident points (\emph{merge points}, tolerance typically
        0.1\,mm) as the first operation of the path-building stage.
        Because point indexing changes after the merge, every cached
        geometric quantity (normals, distance fields) is invalidated at
        that moment --- the class architecture handles this explicitly
        (see \S\ref{sec:impl-object}).
\end{itemize}

% ----------------------------------------------------------------------
\subsection{The distance-to-boundary field and contour shells}
\label{sec:meth-shells}
% ----------------------------------------------------------------------

\subsubsection{Distance field}

Let $\mathcal{S}$ be one layer surface with point set
$P = \{\mathbf{p}_i\}$ and boundary point set
$B \subset P$ (the points belonging to feature/boundary edges). For
every surface point we define
\begin{equation}
  d(\mathbf{p}) \;=\; \min_{\mathbf{b} \in B} \|\mathbf{p} - \mathbf{b}\|_2,
\end{equation}
the geodesic-free (Euclidean) distance to the layer boundary. The field
is computed exactly with a Kd-tree built on $B$: one parallel
nearest-neighbor query ($O(|P| \log |B|)$) returns all distances at
once. The field is stored as the point-data array
\texttt{dist\_to\_boundary} on the surface, so that all subsequent
VTK scalar-field filters operate on it natively.

\subsubsection{Contour shells}

The shells are the isolines of $d$ inside the layer:
\begin{equation}
  \mathcal{C}_c \;=\; \{\,\mathbf{p} \in \mathcal{S} \;:\; d(\mathbf{p}) = c\,\},
  \qquad c \in \{c_1 < c_2 < \dots\}.
\end{equation}
For the piece, the isovalues are multiples of half a nozzle width
(e.g.\ $c \in \{0.5, 1.0, \dots\}$\,mm with a 0.7\,mm nozzle), giving
perimeter passes at successively larger offsets from the layer edge.
The contour filter produces the isolines; a chaining step then
reconstructs ordered polylines from the filter's segment soup
(\S\ref{sec:meth-chaining}).

\subsubsection{Inner zone and offset surfaces}

The infill must not collide with the shells: it is generated only in
the \emph{inner zone}
\begin{equation}
  \mathcal{Z}_{\mathrm{in}} \;=\; \{\,\mathbf{p} \in \mathcal{S} \;:\; d(\mathbf{p}) \ge c_{\max}\,\},
\end{equation}
where $c_{\max}$ is the largest shell isovalue (the \emph{contour
width}). The zone is obtained by scalar clipping: the VTK clip filter
cuts the surface along the isovalue $c_{\max}$ and re-triangulates the
cut cells, so the boundary of $\mathcal{Z}_{\mathrm{in}}$ lies exactly
on the isoline. (Care with conventions is required: in PyVista's
\texttt{clip\_scalar}, the \texttt{invert} flag selects which side is
kept; the slicer documents explicitly that
\texttt{invert=False} keeps $d \ge c_{\max}$.)

% ----------------------------------------------------------------------
\subsection{Fiber direction and fiber-aligned infill planes}
\label{sec:meth-fibers}
% ----------------------------------------------------------------------

\subsubsection{Mean fiber direction by uncentered PCA}

Each layer carries a set of fiber polylines
$F = \{f_i\}$, each $f_i$ being a sequence of 3D points. The local
stiffness of the printed composite is dominated by the fiber
direction, so the matrix infill must be organized \emph{around} the
fibers. The slicer first estimates the dominant fiber direction of the
layer:

\begin{itemize}
  \item collect the unit direction vectors of every fiber segment,
        skipping zero-length segments:
        $\mathbf{u}_{ij} = (f_i^{j+1} - f_i^j)/\|f_i^{j+1} - f_i^j\|$;
  \item build the uncentered second-moment matrix
        $M = \sum_{i,j} \mathbf{u}_{ij}\mathbf{u}_{ij}^\top$;
  \item the mean direction is the dominant eigenvector of $M$ (largest
        eigenvalue).
\end{itemize}

The absence of centering is deliberate and important: a centered PCA
of a bundle of nearly parallel directions would return the axis of
greatest \emph{scatter}, i.e.\ a direction \emph{perpendicular} to the
fibers. The uncentered moment matrix instead has its dominant
eigenvector aligned with the bundle. The sign of the eigenvector is
fixed deterministically (largest-magnitude component positive) so that
the plane generation is reproducible.

\subsubsection{Infill plane family}

Within the layer surface (of mean unit normal $\mathbf{n}$), the
infill direction is set at $\pm 45^\circ$ from the fiber direction
$\mathbf{d}_f$, alternating with the parity of the layer index:
\begin{equation}
  \mathbf{d}_{\pm} \;=\;
  \cos 45^\circ \, \mathbf{d}_f \;\pm\;
  \sin 45^\circ \, (\mathbf{n} \times \mathbf{d}_f),
\end{equation}
normalized. The alternation gives a cross-ply-like stacking of
successive layers, which balances the bending/torsional stiffness of
the laminate. The infill lines themselves are the intersections of
$\mathcal{Z}_{\mathrm{in}}$ with a family of \emph{parallel planes}
\begin{equation}
  \Pi_s \;=\; \{\,\mathbf{p} : (\mathbf{p} - \mathbf{o}) \cdot \mathbf{m} = s\,\},
  \qquad \mathbf{m} = \mathbf{n} \times \mathbf{d}_\pm,
\end{equation}
spaced by the infill spacing $s_{\mathrm{step}}$ (typically 0.5\,mm
for the piece), centered on the projected extent of the inner zone
with a half-step margin so the extreme planes fully bracket the zone.

% ----------------------------------------------------------------------
\subsection{Mold generation: shadow volumes}\label{sec:meth-shadow}
% ----------------------------------------------------------------------

The part must be printed on a support structure ("mold") that
(a) reaches from the printing plane $Z = z_{\mathrm{plane}}$ up to the
part's lower surface, (b) does not intersect the part, and (c) is
itself printable. The slicer constructs the mold as the
\emph{shadow volume} of the part's lower surface with respect to the
build direction. Three constructions are implemented; all operate on
the part's linearized outer surface (the triangulated skin of the
solid mesh) and all produce a volumetric wedge mesh connecting a "top"
point set (on or near the part) to a "bottom" point set (on the
plane or on the part itself).

\subsubsection{Direct projection}

Every point of the part's downward-facing surface
($\mathbf{n}\cdot\mathbf{e}_z < 0$) is projected vertically onto the
plane. The volume is the set of prisms/wedges between each surface
triangle and its vertical projection. This is the fastest method but
builds material even where the part does not overhang (e.g.\ inside
closed cavities facing up).

\subsubsection{Cone of rays}

For each downward-facing triangle, a cone of rays is cast from its
centroid (offset outward along its normal by a small
\texttt{offset\_faces} to avoid self-intersection): one central
vertical ray and \texttt{no\_rays} rays tilted by the half-angle
\texttt{cone\_angle}, uniformly distributed in azimuth
(\texttt{start\_phase} controlling the rotation of the cone). A
(face, ray) pair is \emph{free} if the ray reaches the printing plane
without intersecting the part's surface. For every free pair, the mold
top point is the (offset) centroid and the mold bottom point is the
ray's intersection with the plane --- so the mold's side walls follow
the cone's slope, giving a self-supporting draft angle. The top points
are triangulated by a 2D Delaunay triangulation of their $XY$
projections, and the wedge volume is built on these triangles.

Faces whose centroid projects too close to the part's silhouette (the
convex hull of the $XY$ projection) are discarded
(\texttt{distance\_from\_contour}), which keeps the mold's rim clean.

\subsubsection{Vertical projection with occlusion}

The most physically accurate variant: each \emph{node} of the
lower surface is projected vertically; if the vertical segment from the
node to the plane crosses the part (which happens for nodes below
overhangs, e.g.\ inside the concavity of a C-beam), the bottom point
is the \emph{first intersection} with the part instead of the plane.
The volume therefore fills the gap between the lower surface and the
nearest surface below it (part or plane). A prism-validity test then
discards the triangles whose extrusion would intersect the part: the
segment between the triangle's top centroid and bottom centroid is
ray-cast against the upward-facing part surface, and any real
intersection inside the segment's height range invalidates the prism.

\subsubsection{Wedge volume assembly}

All three methods share the same volume builder: given top points
$\mathbf{t}_i$, bottom points $\mathbf{b}_i$ and triangles
$(a_j, b_j, c_j)$ of the top surface, the volume is the set of
tetrahedralized wedges connecting each triangle
$(\mathbf{t}_{a}, \mathbf{t}_{b}, \mathbf{t}_{c})$ to its projected
triangle $(\mathbf{b}_{a}, \mathbf{b}_{b}, \mathbf{b}_{c})$. The
result is a \texttt{pv.UnstructuredGrid} of tetrahedra, ready for the
same contour-based slicing as the part.

% ----------------------------------------------------------------------
\subsection{Mold slicing and fixed infill grid}\label{sec:meth-moldslice}
% ----------------------------------------------------------------------

The shadow volume is sliced by the simple planar field
\begin{equation}
  q(\mathbf{p}) = \frac{z(\mathbf{p}) - z_{\mathrm{plane}}}{h},
\end{equation}
so that layer $k$ is the horizontal slice at altitude
$z_{\mathrm{plane}} + k\,h$. For each layer, the pipeline computes:

\begin{enumerate}
  \item \textbf{The layer surface}: the isosurface $q = k$, triangulated
        and with point normals.
  \item \textbf{The distance-to-part field}
        $\mathrm{dist\_to\_part}(\mathbf{p}) = \min\!\big(
        \mathrm{dist}(\mathbf{p}, \text{part surface}),
        \ z(\mathbf{p}) - z_{\mathrm{plane}} \big)$.
        The distance to the (triangulated) part surface is obtained
        with an implicit-distance computation; the minimum with the
        height above the plane closes the field inside cavities, where
        the nearest reference is the build plate rather than a far
        wall of the part.
  \item \textbf{Contour shells}: the isolines of
        $\mathrm{dist\_to\_part}$ at the shell isovalues
        $(i + \tfrac12)\, d_{\mathrm{nozzle}}$, giving two shells
        offset from both the part (above the mold) and the free rim
        of the layer.
  \item \textbf{The inner zone}: the clip
        $\mathrm{dist\_to\_part} \ge n_{\mathrm{shells}}\,
        d_{\mathrm{nozzle}}$.
  \item \textbf{The grid infill}: the intersections of the inner zone
        with the two families of planes $x = x_i$ and $y = y_j$.
\end{enumerate}

\paragraph{Fixed infill levels.}
A subtle but decisive detail: the infill plane positions
$\{x_i\}$, $\{y_j\}$ are computed \textbf{once, from the footprint of
layer zero}, and reused for every layer. If each layer used its own
extent, the planes would shift with the shrinking/growing cross-section
of the mold and the beads would not stack vertically --- the printed
wall would be a lattice of misaligned segments instead of continuous
vertical ribs. With fixed levels, each infill line of layer $k$ lies
exactly on top of the corresponding line of layer $k-1$.

\paragraph{Infill spacing.}
For a two-direction grid in which crossings are printed once, the
area fraction covered by beads of width $w$ spaced by $s$ is
$2w/s - (w/s)^2$ (inclusion--exclusion), hence the spacing for a target
fraction $r$:
\begin{equation}
  s \;=\; \frac{w}{1 - \sqrt{1 - r}}.
\end{equation}
A 100\% infill gives $s = w$ (adjacent beads), 20\% gives
$s \approx 4.9\,w$.

% ----------------------------------------------------------------------
\subsection{Segment chaining}\label{sec:meth-chaining}
% ----------------------------------------------------------------------

Both contour filters (isolines of a scalar field) and plane cuts emit
unordered sets of 2-point line segments. The path logic requires
\emph{ordered polylines}. The slicer chains segments with an explicit
graph traversal rather than relying on VTK's stripper, which was
observed to leave multi-level contour output fragmented:

\begin{itemize}
  \item merge coincident points of the segment soup;
  \item compute the degree of every point; if any point has degree
        $> 2$ (a junction), the graph is not a set of simple chains and
        the chaining is abandoned in favor of a slower, more general
        fallback (one slice call per plane);
  \item otherwise build, for each point, its (at most two) neighbors;
  \item walk every open chain from its degree-1 endpoint, and every
        closed loop from any unvisited degree-2 point (the loop is
        closed explicitly by repeating its first point).
\end{itemize}

The complexity is linear in the number of segments, and the walk
naturally preserves the original sampling of the isolines.

% ----------------------------------------------------------------------
\subsection{Infill crossing management (mold)}\label{sec:meth-crossings}
% ----------------------------------------------------------------------

In a grid infill, the primary and secondary families cross each other.
If both were printed fully, every crossing would be deposited twice
(the second pass would melt/cut through the first bead). The mold
algorithm therefore distinguishes, per layer, a \textbf{primary}
direction (X on even layers, Y on odd layers) from a
\textbf{secondary} direction:

\begin{itemize}
  \item primary lines are printed \emph{continuously};
  \item every secondary line is \textbf{cut at each intersection} with
        a primary line, removing a curvilinear zone of one nozzle
        diameter around the crossing point (half a nozzle width on
        each side), so that the secondary beads approach but never
        overlap the primary beads.
\end{itemize}

The intersections are computed in 2D (XY) between the polyline
segments: for a secondary segment $[\mathbf{p}_1, \mathbf{p}_2]$ and a
primary segment $[\mathbf{q}_1, \mathbf{q}_2]$, the parameters
$(t, u)$ of the intersection solve the $2 \times 2$ cross-product
system
\begin{equation}
  \begin{cases}
    \mathbf{p}_1 + t\,\mathbf{d}_1 \equiv \mathbf{q}_1 + u\,\mathbf{d}_2 \quad (\text{in } XY)\\
    \mathbf{d}_1 = \mathbf{p}_2 - \mathbf{p}_1,\quad
    \mathbf{d}_2 = \mathbf{q}_2 - \mathbf{q}_1
  \end{cases}
\end{equation}
and an intersection exists iff $t, u \in [0,1]$ and the segments are
not parallel. The cut positions are converted to curvilinear
abscissae along the secondary polyline, sorted, merged when
overlapping, and the polyline is split into the surviving
sub-polylines, with exact interpolation of the cut endpoints.

A spatial index (Kd-tree on segment midpoints) restricts the pairwise
tests to segments that are close enough to intersect, turning the
naive $O(NM)$ test into an output-sensitive computation.

\subsubsection{Crossing blocker}

Cutting the secondary lines is not sufficient: the \emph{stitching}
stage (\S\ref{sec:meth-stitching}) may connect the end of one line to
the start of the next with a straight printed move that happens to
cross a primary line --- printing the crossing a second time from
above. The mold path builder therefore attaches a
\textbf{connection blocker} to the stitching algorithm: a callable
that tests the candidate connection segment against all primary
segments with the same 2D intersection test, and forbids the
connection when a crossing is detected. A forbidden connection is
replaced by a leave + approach pair (travel moves), exactly like a
too-long connection.

% ----------------------------------------------------------------------
\subsection{Path stitching}\label{sec:meth-stitching}
% ----------------------------------------------------------------------

Each layer yields a set of curves (shells, infill lines) that must be
printed in a single robotic sequence. The slicer stitches them into
one or several \emph{sub-paths} with the following semantics:

\begin{itemize}
  \item a \textbf{curve} is printed;
  \item a \textbf{connection} is a printed straight move from the end
        of one curve to the start of the next, allowed only if it is
        shorter than a maximum length
        (geometrically derived: the diagonal between two adjacent
        infill lines, $\sqrt{s^2 + h^2}$, times a safety margin) and
        not forbidden by a blocker;
  \item an \textbf{approach} is a travel move (not printed) arriving
        at the start of a sub-path, along the local surface normal,
        at distance \texttt{approach\_distance};
  \item a \textbf{leave} is the symmetric travel move leaving the end
        of a sub-path.
\end{itemize}

\subsubsection{Greedy nearest-endpoint algorithm}

The stitching is a greedy loop over an \emph{endpoint table}: for $N$
curves, a $(2N \times 3)$ array holding the first and last point of
each curve, with an activity mask. Starting from a chosen first curve
(optionally reversed), the loop repeatedly:

\begin{enumerate}
  \item computes the (squared, vectorized) distance from the current
        endpoint to all \emph{active} endpoints;
  \item selects the nearest one, identifies its curve and whether the
        nearest point is the curve's end (in which case the curve is
        consumed in reverse);
  \item checks the connection: length and blocker. If legitimate, a
        connection segment is emitted; otherwise the current sub-path is
        closed (leave move) and a new sub-path is opened (approach
        move) at the chosen curve;
  \item consumes the chosen curve, marks its endpoints inactive, and
        repeats until no curve remains.
\end{enumerate}

The final endpoint gets a leave move. Each approach/leave move is
oriented along the \emph{local surface normal} at the corresponding
curve endpoint, so the tool always arrives at and departs from the
surface perpendicularly --- the same direction it will print in.

\subsubsection{Segment codes}

Every segment of the stitched path carries a type and a command code,
kept distinct between the two material systems of the composite part:

\begin{center}
\begin{tabular}{lcc}
\toprule
Move & Pure (matrix / shells) & Composite (fiber / infill) \\
\midrule
Continuation along a curve & 2 & 0 \\
End of a connection & 3 & 1 \\
Approach (start of sub-path) & 5 & 4 \\
\bottomrule
\end{tabular}
\end{center}

The piece path is built as two fully independent passes --- pure
curves (shells) then composite curves (fiber-aligned infill) --- which
are concatenated in the export. This models the physical process:
the matrix and the fiber tows are deposited by different tool states,
and the robot must finish all matrix material of one family before
switching.

% ----------------------------------------------------------------------
\subsection{Tool-pose computation}\label{sec:meth-frames}
% ----------------------------------------------------------------------

A six-DoF robot needs a full orientation at every trajectory point,
not just a position. The slicer defines the tool frame from two
vector fields along the path:

\begin{itemize}
  \item the \textbf{approach normal} $\mathbf{n}(\mathbf{p})$: the unit
        normal of the surface cell closest to $\mathbf{p}$;
  \item the \textbf{travel tangent} $\mathbf{t}(\mathbf{p})$: the unit
        tangent to the trajectory (central difference at interior
        points, one-sided difference at segment endpoints).
\end{itemize}

The orthonormal frame is then built by Gram--Schmidt:
\begin{equation}
  \mathbf{z} = -\mathbf{n} \ (\text{tool axis pointing into the surface}),
  \quad
  \mathbf{x} = \frac{\mathbf{t} - (\mathbf{t}\cdot\mathbf{z})\,\mathbf{z}}
                     {\|\mathbf{t} - (\mathbf{t}\cdot\mathbf{z})\,\mathbf{z}\|},
  \quad
  \mathbf{y} = \mathbf{z} \times \mathbf{x},
\end{equation}
so that $(\mathbf{x}, \mathbf{y}, \mathbf{z})$ is a right-handed basis,
$\mathbf{x}$ the direction of travel and $\mathbf{z}$ the approach
axis. The degenerate case $\mathbf{t} \parallel \mathbf{z}$ (exactly
the case of approach/leave moves, which travel along the normal) is
handled by an arbitrary perpendicular fallback, since travel
direction carries no meaningful orientation information there.

Two aspects deserve emphasis. First, the tool is \emph{not}
rotationally symmetric about its approach axis: the choice of
$\mathbf{x}$ fixes the tool's rotation about its own axis, and getting
it wrong can mean the tool body colliding with the part even when the
approach axis is correct. Second, the \textbf{mold redefines the
approach normal}: for support material, the natural approach direction
is perpendicular to the \emph{part surface} the support is printed
against, not to the support layer's own surface. The mold class
therefore overrides the normal lookup to sample the part's surface
normals instead --- one of the two places where the class hierarchy
changes behavior rather than merely adding it.

Tangents are computed independently within each segment, so a small
discontinuity in the frame's $\mathbf{x}$ axis can occur exactly at a
segment boundary (curve $\to$ connection): only the rotation of the
tool about its approach axis is affected, never the position, and the
robot interpolates smoothly over it.

% ----------------------------------------------------------------------
\subsection{Polyline simplification}\label{sec:meth-simplify}
% ----------------------------------------------------------------------

Raw isolines can contain thousands of nearly collinear points. The
export stage offers two levels of simplification:

\begin{itemize}
  \item \textbf{Ramer--Douglas--Peucker (RDP)} \cite{rdp}, implemented
        iteratively (explicit stack, no recursion) and directly in 3D:
        the perpendicular distance of each candidate point to the
        current 3D chord is computed exactly; points are kept as long
        as the maximum deviation exceeds the tolerance. Endpoints are
        always kept, and simplification is applied per segment, so
        distinct segments are never merged.
  \item \textbf{Infill collapse} (mold): a mold infill line is, in the
        normal case, the intersection of a planar layer with a plane
        $x = \text{const}$ or $y = \text{const}$ --- a straight line.
        It can therefore be reduced to its two endpoints, shrinking
        the trajectory file by orders of magnitude. A safety net
        guards this collapse: the line is flattened only if every
        point lies within a tolerance of the endpoint chord \emph{and}
        the points progress monotonically along the chord. A line that
        doubles back (contains a printed jump) is left untouched, and
        a merely curved line falls back to RDP; no gap is ever erased.
\end{itemize}

% ----------------------------------------------------------------------
\subsection{Trajectory export}\label{sec:meth-export}
% ----------------------------------------------------------------------

The final trajectory is a plain-text table, one row per trajectory
point:
\begin{center}
\texttt{x y z ux uy uz vx vy vz wx wy wz type [layer]}
\end{center}
where $(u\!\cdot\!\cdot, v\!\cdot\!\cdot, w\!\cdot\!\cdot)$ are the
three columns of the rotation matrix $R$ (the tool's X, Y and Z axes
in world coordinates), \texttt{type} is the segment label
(\texttt{approach}, \texttt{curve}, \texttt{connection},
\texttt{leave}) and \texttt{layer} the layer index. A small
post-processing drops the duplicated point at each segment junction
(consecutive segments share one point). Multi-layer export writes the
header once and appends every subsequent layer to the same file, in
print order.

This format is deliberately minimal and self-describing: building the
$4\times 4$ pose matrix on the receiving end is a reshape, and the
type labels let the robot-side post-processor map moves to printed /
travel motions and select the deposition tool state.
